\chapter{Конечный вихрь сцеплённой фазы}

Совместный вакуум предыдущей главы равен $SO(3)/Z_3$. Поэтому
\begin{equation}
 \pi_1(SO(3)/Z_3)=Z_6,
 \qquad \pi_2(SO(3)/Z_3)=0.
\end{equation}
Прежний проекторный ёж не является изолированным зарядом полного
вакуума. Минимальный устойчивый сектор теперь имеет коразмерность два.

Выберем генератор $J$ вращений вокруг выбранной тетраэдрической оси и
разложим $T_*=T_0+T_3$ на инвариантную и вращающуюся части. Используется
анзац
\begin{align}
 Q&=Q_*,\\
 T&=b(r)T_0+a(r)\rho(e^{\theta J/3})T_3,\\
 B&=-\frac13K(r)J\,d\theta.
\end{align}
После полного обхода возникает элемент $Z_3$, а его подъём в $SU(2)$
имеет порядок шесть. Граничные условия равны
\begin{equation}
 K(0)=a(0)=0,quad b'(0)=0,qquad
 K(\infty)=a(\infty)=b(\infty)=1.
\end{equation}

Ограничение уже выведенного потенциала имеет вид
\begin{equation}
\begin{split}
V={}&\frac{31744}{19683}-\frac{38912}{19683}b^2
+\frac{5120}{6561}b^4-\frac{8192}{6561}a^2\\
&+\frac{8192}{19683}a^2b^2+\frac{8192}{19683}a^4.
\end{split}
\end{equation}
Нормированный триплетный след задаёт коэффициенты
$A=C=128/81$, $B=160/81$, $G=2/27$ без новых весов.

Краевая задача решена на $10^{-5}\le r\le20$. Получено
\begin{equation}
 b(0)=1.0708503824\ldots,qquad
 a'(0)=2.562227029\ldots,qquad
 \mathcal T=3.1060360712\ldots.
\end{equation}
Максимальный относительный остаток равен $2.0\cdot10^{-7}$, относительный
вириальный остаток --- $4.96\cdot10^{-8}$.

\begin{protocol}
\begin{enumerate}
 \item топологический класс: \textbf{$Z_6$};
 \item гладкий профиль: \textbf{получен};
 \item безразмерное натяжение: \textbf{конечно};
 \item абсолютный масштаб: \textbf{не выведен};
 \item полная устойчивость: \textbf{не проверена};
 \item следующий гейт: \textbf{радиальный гессиан вихря}.
\end{enumerate}
\end{protocol}

Машинный сертификат сохранён в
\path{s2t_v6_bosonic_defect_q_tetrahedral_coupled_defect_profile_gate_results.json}.